% ==============================================================================
% US-072 -- FarSLIP contrastivo-fenológico (sección Method/Experiments modular).
%
% Sección auto-contenida del Paper Track (E11). Documenta el diagnóstico del bug
% torch.randn, el fix con prototipos fenológicos reales, el filtro 3:1 Meadow
% per-patch, la ablación de bandas, el protocolo incremental 4->18 y el marco
% Be My Eyes. TODAS las cifras provienen de artefactos REALES del repo (ver los
% comentarios de cada \input{...}). Las tablas viven en
% paper/tables/farslip-method/ y se incluyen al final de cada subseccion.
%
% Para compilar de forma autonoma (fuera del manuscrito principal) se requiere:
%   \graphicspath{{../figures/}{../../reports/farslip/figures/}}
% y los \bibitem li2025farslip / wen2025phenology / huang2025bemyeyes
% (ver paper/bib/farslip_refs.bib). Si se incrusta en
% docs/final_doc/Avance7_equipo17.tex, esos recursos ya existen alli.
% ==============================================================================

\section{FarSLIP contrastivo-fenológico: del fallo contra ruido a la señal
visión-lenguaje}
\label{sec:farslip_method}

Esta sección documenta el camino metodológico que eleva FarSLIP de una
ablación negativa --- perdía F1-macro $0.163$ frente al $0.233$ de AlphaEarth
sobre las 18 clases de PASTIS-R --- a un componente visión-lenguaje cuyo aporte
es \textit{complementario y honesto}. El relato es deliberadamente transparente:
reportamos el bug que originaba el resultado negativo, el fix que cierra el gap
\textit{propio} de FarSLIP, y el hecho de que, aun corregido, el espacio FarSLIP
fiel \textbf{no supera} a AlphaEarth como clasificador denso. Su valor reside en
el ensamble heterogéneo y en la transferibilidad multirregión (Sec.~\ref{sec:farslip_transfer}).

% ------------------------------------------------------------------------------
\subsection{F1 --- Diagnóstico del bug: alineamiento contrastivo contra ruido}
\label{sec:farslip_f1}

La pérdida de alineamiento región-categoria de FarSLIP~\cite{li2025farslip}
(\texttt{RegionCategoryAlignmentLoss}, un InfoNCE con temperatura $0.07$) alinea
el token \texttt{[CLS]} del estudiante contra una matriz de \textit{prototipos
de texto}. En la implementación original
(\texttt{ml/farslip/train.py}, ruta \texttt{proto\_source="random"}) esa matriz
se inicializaba con \texttt{torch.randn(n\_protos, hidden\_dim)}: ruido gaussiano
sin contenido semántico. El contraste empujaba el \texttt{[CLS]} hacia vectores
aleatorios, de modo que la señal categórica nunca se anclaba a una descripción
real del cultivo. El docstring de \texttt{build\_text\_prototypes} nombra el
síntoma textualmente: \textit{``contrastive alignment against noise --- the bug
that made FarSLIP lose 0.163 vs 0.233''}. Este es el origen del resultado
negativo: no un fallo del modelo de fundación, sino una inicialización que
convertía una pérdida contrastiva en regularización hacia ruido.

% ------------------------------------------------------------------------------
\subsection{F2 --- Prototipos fenológicos reales y reproyección ortogonal frozen}
\label{sec:farslip_f2}

El fix sustituye el ruido por \textit{prototipos fenológicos reales}, una clase
a la vez. Para cada una de las 18 clases de PASTIS-R se redacta una descripción
fenológica siguiendo el marco de Wen et al.~\cite{wen2025phenology} (inicio de
verdor, pico de NDVI, fecha de senescencia, textura de dosel), generada con
Gemini~Flash y codificada con MiniLM (\texttt{all-MiniLM-L6-v2}) a un embedding
de 384 dimensiones L2-normalizado. El flujo (ruta \texttt{proto\_source="pastis"})
es:

\begin{enumerate}[leftmargin=*]
    \item \texttt{load\_class\_prototype\_embeddings()} $\rightarrow$
    \texttt{proto\_18} $(18, 384)$, prototipos MiniLM por clase.
    \item \texttt{expand\_to\_cap()} $\rightarrow$ \texttt{proto\_cap}
    $(32, 384)$, mapeo del esquema PASTIS-18 al vocabulario CAP de 32 categorias.
    \item \texttt{np.tile()} $\rightarrow$ \texttt{proto\_tiled}
    $(n\_regions \cdot 32, 384)$ en orden \textit{región-major}
    (\texttt{región $\cdot$ n\_categories + category}), coincidente con el target
    de la pérdida.
\end{enumerate}

El \textit{tile} de 384 dimensiones se reproyecta a la dimensión del
\texttt{[CLS]} del estudiante mediante una matriz semi-ortogonal $W$ generada
con \texttt{torch.nn.init.orthogonal\_}, \texttt{requires\_grad=False}: nunca
entra al optimizador ni al \texttt{state\_dict}. Al ser ortogonal preserva la
norma $L_2$ y los ángulos del subespacio MiniLM, de modo que las 18 clases
permanecen distinguibles tras la reproyección.

\paragraph{Corrección factual sobre la dimensionalidad.}
La reproyección de los prototipos del \textit{loss} es \textbf{$384\rightarrow768$},
no $384\rightarrow512$: el \texttt{hidden\_size} del teacher CLIP ViT-B/16 es
$768$, confirmado por \texttt{trainer.teacher.config.hidden\_size} y por la
columna \texttt{n\_dims}$=768$ del artefacto de evaluación. La cifra ``512-dim
por parcela'' corresponde a otro objeto --- el banco de \textit{features} que el
extractor produce a nivel de parcela --- y no debe confundirse con la
reproyección de prototipos del InfoNCE. Distinguimos explícitamente:
banco-parcela $512$-dim frente a prototipos/\texttt{[CLS]} $768$-dim.

\paragraph{Efecto del fix sobre el gap propio.}
El fix cierra el gap \textit{propio} de FarSLIP: el espacio fiel v2 pasa de
F1-macro $0.163$ (contra ruido) a $0.5551 \pm 0.0205$ (con prototipos
fenológicos) sobre las mismas 567 parcelas, una mejora de $+0.392$ respecto al
baseline buggeado. La Tabla~\ref{tab:farslip_vs_alphaearth} reporta la
comparación honesta frente a AlphaEarth.

\input{../tables/farslip-method/farslip_vs_alphaearth.tex}

\paragraph{Claim defendible.}
Corregido el bug, el espacio FarSLIP fiel \textbf{sigue por debajo} de AlphaEarth
($-0.0024$ en F1-macro, $-0.0025$ en silhouette). No afirmamos que FarSLIP supere
a AlphaEarth ni a Gemini: el aporte de FarSLIP es como miembro complementario del
ensamble y en el transfer multirregión, no como clasificador denso ganador. Esta
es la lectura que el resto de la sección sostiene.

% ------------------------------------------------------------------------------
\subsection{F3 --- Filtro 3:1 Meadow por parcela (\texttt{dominance\_ratio})}
\label{sec:farslip_f3}

PASTIS-R está dominado por \textit{Meadow} (clase 1) y \textit{Background}
(clase 0), que juntos pueden cubrir $>80\%$ de un patch. Extraer
\textit{features} sobre esos patches contamina el espacio de embeddings con
señal poco informativa. El filtro previo de Isaac operaba en modo
\texttt{coverage} (legacy): conservaba un patch si la cobertura combinada de las
clases objetivo sobre los pixeles validos superaba un umbral global. El presente
trabajo introduce el modo \texttt{dominance\_ratio} en
\texttt{ml/data/pastis\_filter.py}: conserva un patch si
$\text{Meadow}_{px} \leq 3 \cdot \text{second}_{px}$, donde
$\text{second}_{px}$ es el conteo mayor \textit{en ese mismo patch} entre las
clases objetivo, excluyendo Meadow, Background~(0) y Void~(19) por identificador
literal (no vía \texttt{ignore\_index}) para no contaminar la razon. La
comparación es inclusiva ($\leq$): un patch exactamente 3:1 se conserva.

La diferencia con el filtro \texttt{coverage} es conceptual: \texttt{coverage}
mide pureza \textit{global} de clases objetivo; \texttt{dominance\_ratio} acota
el desbalance \textit{local} de Meadow respecto a la segunda clase de cada patch,
descartando solo los patches sobre-dominados por prado sin tirar la señal
minoritaria del resto. Casos borde explícitos y testeados:
$\text{Meadow}_{px}=0$ conserva (métrica $0.0$); $\text{Meadow}_{px}>0$ con
$\text{second}_{px}=0$ (sin clase competidora) descarta (métrica $\infty$).

% ------------------------------------------------------------------------------
\subsection{F4 --- Ablación de bandas (rgb / nir\_rgb / 4band)}
\label{sec:farslip_f4}

Se evaluaron tres variantes de seleccion de bandas (US-035), implementadas como
\textit{transform} por recorte en \texttt{ml/farslip/bands.py}: \textit{rgb}
(B04-B03-B02, 3 canales), \textit{nir\_rgb} (B08-B04-B03 falso color, 3 canales)
y \textit{4band} (B02-B03-B04-B08 identidad, 4 canales). La variante 4band
reutiliza la adaptación $3\rightarrow4$ del \texttt{patch\_embed} con
NIR\,=\,mean(RGB) (US-034) para evitar neuronas muertas al ampliar el primer
bloque convolucional.

\input{../tables/farslip-method/band_ablation.tex}

\paragraph{Caveat de evaluación (resultado parcial honesto).}
Las tres corridas existen y convergen en pérdida de destilación
(Tabla~\ref{tab:band_ablation}): la pérdida total final es muy similar entre
\textit{rgb} ($2.2893$) y \textit{nir\_rgb} ($2.2858$), y algo mayor en
\textit{4band} ($2.3553$), con una pérdida de \textit{patch} notablemente más
alta en 4band ($0.0786$ frente a $\approx0.018$) coherente con el canal NIR
adicional. Sin embargo, los logs no registran la métrica de evaluación
consolidada (F1-macro/mIoU) por variante; por tanto NO reportamos un ganador de
clasificación entre bandas. Reportar una cifra de F1 inexistente sería inventar
datos. El cierre de esta tabla con la métrica de evaluación queda como tarea
pendiente documentada en los \textit{blockers} de la épica.

% ------------------------------------------------------------------------------
\subsection{F5 --- Protocolo incremental 4$\rightarrow$18 y fusion en el ensamble}
\label{sec:farslip_f5}

El clasificador FarSLIP se construye de forma incremental
(\texttt{ml/farslip/incremental\_curriculum.py}): se parte de las 4 clases más
cardinales (Stage-1) y se añaden 2 clases por paso hasta 18 (Stage-2), inicializando
cada etapa desde la anterior con \texttt{strict=False} para permitir el cambio de
dimensión del cabezal. El criterio de parada (\texttt{stop\_criterion}) contempla
tres razones explícitas: \textit{(i)} clases nuevas inaceptables
(F1$<0.30$), \textit{(ii)} degradación de clases previas
($\text{drop}>0.05$, olvido catastrófico), \textit{(iii)} máximo de clases
alcanzado.

La Tabla~\ref{tab:cardinality_sweep} muestra el barrido de cardinalidad por
parcela --- la evidencia empírica que motiva el protocolo. La dificultad crece
monotónicamente con el número de clases: F1-macro $0.7025$ con 4 clases,
$0.4579$ con 6, $0.4075$ con 8, $0.3589$ con 10 y $0.3328$ con 12. El número de
parcelas evaluables crece (de 1301 a 3200) al admitir clases más raras, pero la
calidad promedio cae: las clases minoritarias con fenología ambigua arrastran el
F1-macro.

\input{../tables/farslip-method/cardinality_sweep.tex}

\paragraph{Fallback honesto a 18 clases desde cero.}
El baseline fiel de 18 clases sin pesos de clase (variante
\texttt{faithful\_v2}, \textit{región\_category} MPCL + $L_{glo}$) alcanza
F1-macro $0.164$ y mIoU $0.111$, con 4 clases de señal clara
(Meadow $0.761$, Orchard $0.516$, Grapevine $0.489$, Corn $0.464$) y 14 clases
en cero o muy bajo. La variante con pesos de clase \textit{inverse-frequency}
(\texttt{faithful\_v2\_cw}) EMPEORA el balance neto a F1-macro $0.102$: desinfla
los dominantes más de lo que rescata a los raros. Estos resultados se reportan
sin maquillaje; las 18 clases son $\approx4.5\times$ más difíciles que las 4
fáciles ($0.576$ en US-036-a, no comparable directamente).

\paragraph{Fusion en el ensamble heterogéneo.}
El aporte neto de FarSLIP se mide en el meta-modelo de \textit{stacking}. El
ensamble de 5 miembros (TSViT-pheno, U-TAE, XGB-AlphaEarth, FarSLIP-ft18,
FarSLIP-zeroshot) sobre probabilidades \textit{out-of-fold} alcanza F1-macro
$0.6477$ (\textit{accuracy} $0.7935$) frente a $0.6359$ del ensamble de 3
miembros sin FarSLIP: un delta de $+0.0118$ atribuible a las dos ramas FarSLIP en
el \textit{stacking} OOF. El manuscrito principal reporta la cifra del ensamble
final consolidado (F1-macro $0.749$); la magnitud exacta de la mejora atribuible
a FarSLIP en esa cifra final difiere de este delta OOF y se trata como decision
editorial (ver \textit{blockers}). En todo caso, la conclusion es estable: el
aporte de FarSLIP es \textbf{positivo pero pequeño} ($\approx+0.3\%$ a $+1.2\%$
según la métrica), nunca el motor del clasificador.

% ------------------------------------------------------------------------------
\subsection{Marco \textit{Be My Eyes}: el LLM no clasifica pixeles}
\label{sec:farslip_bemyeyes}

El resultado de que los prototipos fenológicos de texto aporten solo
$\approx0.3\%$ como clasificador supervisado --- y no el $5\%$ que una lectura
optimista sugeriría --- es un resultado \textit{esperado y correctamente
enmarcado}, no un fracaso. Bajo el patron \textit{Be My Eyes} de Huang et
al.~\cite{huang2025bemyeyes}, el sistema separa \textit{perceiver} y
\textit{reasoner}: los \textit{perceivers} son nuestros modelos densos
(TSViT-pheno, U-TAE) y FarSLIP, que \textit{ven} cada parcela y emiten
observaciones; el \textit{reasoner} es un LLM frontera \textit{frozen}
(Gemini~3.5~Flash en la nube, o Qwen3-30B-A3B on-prem para soberania de
datos) que razona sobre esas observaciones en texto plano. El LLM \textbf{no es
un clasificador de pixeles}; el clasificador final es el ensamble heterogéneo.

Por tanto, la descripción fenológica de FarSLIP no debe juzgarse como cabezal de
clasificación supervisada (donde su aporte es marginal), sino como \textit{señal
de anclaje} que enriquece la explicacion del \textit{reasoner} sin enviar
tensores ni logits al LLM y sin añadir latencia perceptible. Esta separación
explica por que el resultado negativo de F2 como clasificador es compatible con
el valor positivo de FarSLIP como componente visión-lenguaje del copiloto.

% ------------------------------------------------------------------------------
\subsection{Caveats explícitos}
\label{sec:farslip_caveats}

\begin{itemize}[leftmargin=*]
    \item \textbf{Reproyección cruda.} El lift ortogonal $384\rightarrow768$ es
    un mapeo lineal que preserva norma y ángulos, pero \textit{no} es el
    \texttt{CLIPTextModel} nativo: es una aproximación (caveat
    \texttt{ortho\_proj\_crude\_approx} registrado en los parámetros de MLflow).
    La ruta con el codificador de texto nativo queda como trabajo posterior.
    \item \textbf{Tamaño de muestra.} La evaluación FarSLIP-vs-AlphaEarth se
    realiza sobre $n=567$ parcelas (orden $\approx500$--$700$ por fold), con
    validación cruzada espacial. Las cifras llevan su desviación estándar.
    \item \textbf{Pureza de parcela.} El filtro \texttt{dominance\_ratio} y la
    agregación a nivel de parcela mantienen una pureza alta ($\approx98\%$) tras
    descartar los patches sobre-dominados por Meadow; el residual de mezcla se
    asume en la etiqueta de parcela dominante.
    \item \textbf{Saturación del benchmark.} FarSLIP fiel no supera a AlphaEarth
    como clasificador denso ($-0.0024$); su valor es complementario (ensamble) y
    de transferibilidad, coherente con la transferibilidad espacial limitada
    reportada para embeddings tipo AlphaEarth.
\end{itemize}
